% Revize 2026-10-27: cílené doplnění; nové úlohy a řešení jsou vlastní tvorba.
% Metodický vzor: 4_rocnik/pripravy/Uvod_do_kombinatoriky.pdf (samostatný dokument).
% Hodina 38, 1. ročník SPŠ. Rozsah: Rozpis_vyuky_MAT_1_rocnik.tex, řádek 38.

% Plán: 1_rocnik/tematicke_plany/SPS.pdf; společná projektová šablona.

% Učebnice: ucebnice/01_Zakladni_poznatky/Mocniny_a_odmocniny/Mocniny a Odmocniny.tex

% Původ úloh: Vlastní zadání a vlastní formulace výkladu; metodická inspirace zdroji.

% Metodický zdroj: externi_zdroje/realisticky/Mocniny/08_Exponencialni_tvar_cisla.pdf

% URL: http://www.realisticky.cz/ucebnice/01%20Matematika%20S%C5%A0/01%20Z%C3%A1kladn%C3%AD%20poznatky/06%20Mocniny/08%20Exponenci%C3%A1ln%C3%AD%20tvar%20%C4%8D%C3%ADsla.pdf

% Metodický zdroj: externi_zdroje/realisticky/Odmocniny/07_Druha_odmocnina.pdf

% URL: http://www.realisticky.cz/ucebnice/01%20Matematika%20S%C5%A0/01%20Z%C3%A1kladn%C3%AD%20poznatky/02%20%C4%8C%C3%ADseln%C3%A9%20obory/07%20Druh%C3%A1%20odmocnina.pdf

% Metodický zdroj: externi_zdroje/realisticky/Odmocniny/08_Treti_odmocnina.pdf

% URL: http://www.realisticky.cz/ucebnice/01%20Matematika%20S%C5%A0/01%20Z%C3%A1kladn%C3%AD%20poznatky/02%20%C4%8C%C3%ADseln%C3%A9%20obory/08%20T%C5%99et%C3%AD%20odmocnina.pdf

\documentclass[11pt,a4paper]{article}

% Spolecny styl pro ucitelske pripravy MAT (Work/MAT)
% Verze 1.1, 2026-10-23
\NeedsTeXFormat{LaTeX2e}
\ProvidesPackage{priprava}[2026/10/23 v1.1 Sablona ucitelskych priprav MAT]

% Zakladni typografie a cestina
\RequirePackage[T1]{fontenc}
\RequirePackage[utf8]{inputenc}
\RequirePackage[czech]{babel}
\RequirePackage{lmodern}
\RequirePackage{microtype}
\RequirePackage{amsmath,amssymb}
\RequirePackage{xcolor}
\RequirePackage{enumitem}
\RequirePackage[most]{tcolorbox}
\RequirePackage{tabularx,booktabs,array}
\RequirePackage{fancyhdr}
\RequirePackage{needspace}
\RequirePackage{geometry}
\RequirePackage[hidelinks]{hyperref}

\geometry{a4paper,top=17mm,bottom=18mm,left=23mm,right=23mm}
\definecolor{notegray}{gray}{0.39}
\setlength{\parindent}{0pt}
\setlength{\parskip}{2.5pt plus 1pt minus 0.5pt}
\setlist[itemize]{topsep=2pt,itemsep=1pt,parsep=0pt,leftmargin=1.45em}
\setlist[enumerate]{topsep=2pt,itemsep=2pt,parsep=0pt,leftmargin=1.65em}
\setcounter{secnumdepth}{0}
\renewcommand{\arraystretch}{1.12}

% Zahlavi: vlevo rocnik | globalni tema, vpravo cislo: nazev hodiny
% Zapati zustava zcela prazdne.
\pagestyle{fancy}
\fancyhf{}
\renewcommand{\headrulewidth}{0.35pt}
\renewcommand{\footrulewidth}{0pt}
\setlength{\headheight}{14pt}
\newcommand{\nastavPripravu}[4]{%
  \fancyhead[L]{\small #1\enspace\textbar\enspace #2}%
  \fancyhead[R]{\small #3: #4}%
}

% Minimalisticke cernobile boxy. Kazda definice zvlast.
\tcbset{commonlesson/.style={enhanced,breakable=false,colback=white,colframe=black,
  colbacktitle=white,coltitle=black,boxrule=0.45pt,
  arc=0.8mm,outer arc=0.8mm,left=3mm,right=3mm,top=1.6mm,bottom=1.6mm,
  toptitle=1mm,bottomtitle=0.7mm,fonttitle=\bfseries,
  attach boxed title to top left={xshift=2.5mm,yshift=-1.15mm},
  boxed title style={colback=white,colframe=white,boxrule=0pt,left=1mm,right=1mm}}}
\newtcolorbox{zadaniBox}[1][]{commonlesson,title={Zadání},#1}
\newtcolorbox{prikladBox}[1][]{commonlesson,title={Příklad},#1}
\newtcolorbox{definiceBox}[1][]{commonlesson,title={Definice},#1}

% Spolecny box pro zadani a jeho reseni, bez mezery mezi castmi.
% Pouziti:
%   \begin{prikladReseniBox}
%     Zadani ulohy.
%     \tcblower
%     \begin{align*} ... \end{align*}
%   \end{prikladReseniBox}
% Spodni cast ma jemny sedy podklad a delici linku; lze v ni sazet
% vice radku matematiky. Obe casti maji jen jeden spolecny ramecek.
\tcbset{solutionlesson/.style={
  commonlesson,
  enhanced jigsaw,
  breakable,
  colbacklower=black!3,
  segmentation style={solid,draw=black!25,line width=0.35pt},
  fontlower=\small\color{notegray},
  before lower={\noindent\textbf{Řešení:}\enspace},
  leftlower=3mm,rightlower=3mm
}}
\newtcolorbox{prikladReseniBox}[1][]{solutionlesson,title={Příklad},#1}
\newtcolorbox{zadaniReseniBox}[1][]{solutionlesson,title={Zadání},#1}

% Pomucky pro autora pripravy: sedive poznamky; \reseni pro starsi pripravy
\newcommand{\poznamka}[1]{\par\begingroup\small\color{notegray}\itshape
  \textnormal{\textbf{Poznámka pro učitele.}}\enspace #1\par\endgroup}
\newcommand{\reseni}[1]{\par\begingroup\small\color{notegray}
  \textbf{Očekávané řešení:}\enspace #1\par\endgroup}
\newcommand{\blok}[1]{\needspace{5\baselineskip}\subsection*{#1}}
\newcommand{\vyrok}[1]{\uv{#1}}


\nastavPripravu{1. ročník}{Odmocniny}{38}{Slovní a technické úlohy}

\begin{document}

\blok{Výklad a souvislosti}

V praktické úloze nejprve rozhodni, zda je známá délka, obsah, nebo objem. Teprve potom zvol mocninu nebo odmocninu. Výsledek obsahuje jednotku a musí odpovídat velikosti původních údajů.

\begin{definiceBox}[title={Připomenutí definice},breakable=false]
\textbf{Druhá odmocnina} z obsahu čtverce $S\ge0$ určuje jeho nezápornou délku strany $a=\sqrt S$, protože $a^2=S$. Číselné údaje musí být vyjádřeny v odpovídajících jednotkách.
\end{definiceBox}


Pro čtverec použijeme $S=a^2$, pro krychli $V=a^3$. Délky měříme v jednotkách délky, obsahy v jejich druhých mocninách a objemy ve třetích mocninách. Při převodu $1\,\mathrm m=10^3\,\mathrm{mm}$ proto platí $1\,\mathrm m^2=10^6\,\mathrm{mm}^2$. V modelu vyber fyzikálně smysluplnou nezápornou délku.

\blok{Řešené příklady}

\begin{prikladReseniBox}[title={Příklad 1},breakable=false]
Čtvercová deska má obsah $72\,\mathrm{cm}^2$. Urči přesně stranu a úhlopříčku.
\tcblower
Strana je $a=\sqrt{72}\,\mathrm{cm}=6\sqrt2\,\mathrm{cm}$. Z Pythagorovy věty je úhlopříčka $u=\sqrt{72+72}\,\mathrm{cm}=12\,\mathrm{cm}$. Kontrola obsahu: $(6\sqrt2)^2=72$.
\end{prikladReseniBox}


\begin{prikladReseniBox}[title={Příklad 2},breakable=false]
Objem krychlového dílu je $54\,\mathrm{cm}^3$. Urči hranu v přesném tvaru.
\tcblower
$a=\sqrt[3]{54}\,\mathrm{cm}=3\sqrt[3]2\,\mathrm{cm}$. Kontrola: $(3\sqrt[3]2)^3=27\cdot2=54$. Z $27<54<64$ plyne $3\,\mathrm{cm}<a<4\,\mathrm{cm}$.
\end{prikladReseniBox}


\blok{Samostatné procvičení}

\begin{zadaniReseniBox}[title={Procvičení},breakable=false]
\begin{enumerate}[label=\alph*)]
\item Urči stranu čtverce s obsahem $50\,\mathrm{mm}^2$.
\item Urči hranu krychle o objemu $0{,}008\,\mathrm m^3$.
\item Na $2\cdot10^3$ dílů připadá po $3\cdot10^{-4}\,\mathrm m^2$ materiálu. Jaký je celkový obsah bez odpadu?
\end{enumerate}
\tcblower
\begin{enumerate}[label=\alph*)]
\item $a=\sqrt{50}\,\mathrm{mm}=5\sqrt2\,\mathrm{mm}$. Druhá mocnina číselné hodnoty je $50$.
\item $a=0{,}2\,\mathrm m$, neboť $0{,}2^3=0{,}008$.
\item $(2\cdot10^3)(3\cdot10^{-4})\,\mathrm m^2=6\cdot10^{-1}\,\mathrm m^2=0{,}6\,\mathrm m^2$.
\end{enumerate}
\end{zadaniReseniBox}


\poznamka{Nehodnoť pouze výpočet: žák musí zdůvodnit volbu odmocniny a jednotky. Nezaváděj nové fyzikální vztahy. Při nesprávném převodu obsahu nech žáka násobit dvě převedené délky; samotné přepsání předpony nestačí.}


\blok{Navazující procvičování ve třídě}

\begin{zadaniReseniBox}[title={Rezerva: škálování a kontrola jednotek},breakable=false]
\begin{enumerate}[label=\alph*)]
\item Hranu krychle zvětšíme třikrát. Kolikrát se zvětší objem? Zdůvodni bez dosazení konkrétní délky.
\item Čtvercová fólie má obsah $0{,}0049\,\mathrm m^2$. Urči stranu v centimetrech.
\end{enumerate}
\tcblower
\begin{enumerate}[label=\alph*)]
\item Z $V=a^3$ a $V_{\mathrm{nový}}=(3a)^3=27a^3$ plyne pro $a>0$ zvětšení $27$krát.
\item $a=\sqrt{0{,}0049}\,\mathrm m=0{,}07\,\mathrm m=7\,\mathrm{cm}$.
\end{enumerate}
\end{zadaniReseniBox}

\blok{Volitelná rezerva}

\begin{prikladReseniBox}[title={Rozšíření procvičení},breakable=false]
Čtverec má stranu $2\cdot10^{-3}\,\mathrm m$. Urči obsah v $\mathrm m^2$ a $\mathrm{mm}^2$.
\tcblower
$S=4\cdot10^{-6}\,\mathrm m^2=4\,\mathrm{mm}^2$. Nezávisle: strana je $2\,\mathrm{mm}$, takže obsah je $4\,\mathrm{mm}^2$.
\end{prikladReseniBox}



\poznamka{Základ: původní slovní úlohy, u nichž žák vysvětluje model a jednotky. Nová dvojice je rezerva, nikoli nová fyzika. Domácí čtveřice umožňuje nezávislou zkoušku umocněním a převodem délky před výpočtem obsahu. Rychlí žáci mají porovnat oba kontrolní postupy.}

\blok{Domácí procvičení}

\begin{zadaniReseniBox}[title={Domácí úlohy},breakable=false]
\begin{enumerate}[label=\alph*)]
\item Urči stranu čtverce s obsahem $98\,\mathrm{cm}^2$.
\item Urči hranu krychle s objemem $0{,}000125\,\mathrm m^3$ v centimetrech.
\item Každý z $4\cdot10^3$ dílů má objem $2{,}5\cdot10^{-6}\,\mathrm m^3$. Urči jejich celkový objem.
\item Čtverec má stranu $3\cdot10^{-2}\,\mathrm m$. Urči obsah v $\mathrm m^2$ a $\mathrm{cm}^2$.
\end{enumerate}
\tcblower
\begin{enumerate}[label=\alph*)]
\item $7\sqrt2\,\mathrm{cm}$.
\item $0{,}05\,\mathrm m=5\,\mathrm{cm}$; kontrola $0{,}05^3=0{,}000125$.
\item $10\cdot10^{-3}\,\mathrm m^3=0{,}01\,\mathrm m^3$.
\item $9\cdot10^{-4}\,\mathrm m^2=9\,\mathrm{cm}^2$.
\end{enumerate}
\end{zadaniReseniBox}

\end{document}
